\section{Introduction}

When a language model performs well on MMLU and poorly on HellaSwag, practitioners instinctively ask: what is different about these benchmarks? The natural follow-up question is almost never asked: what is different about the \textit{data} that produced this asymmetry? Data mixing decisions during pre-training --- which domains to include, in what proportions, across how many tokens --- are among the most impactful and least understood levers in large language model development. We set out to measure which domains of The Pile \cite{Gao2020Pile} drive which NLP benchmarks --- and found that the infrastructure to answer this question has existed for years, the domains are measurably distinct with very large effect sizes, but the causal pathway from domain exposure to benchmark-specific capability improvement is murkier than the field has assumed.

The problem has three nested layers. At the surface, it is well established that data composition matters: DCLM \cite{Li2024DCLM} shows that model-based filtering outperforms heuristic filtering (64\% vs.\ $\sim$57\% MMLU at 7B), and domain mixing frameworks like RegMix \cite{Liu2024RegMix} and DoReMi \cite{DoReMi2023} demonstrate that optimal domain weights improve aggregate validation loss. But these analyses report \textit{aggregate} effects --- a single number summarizing performance across dozens of tasks. Beneath this surface lies a deeper gap: no study has produced a per-domain $\times$ per-benchmark coefficient matrix from a controlled within-family analysis. We do not know, empirically, whether Wikipedia exposure specifically improves MMLU more than HellaSwag, or whether Books exposure specifically improves commonsense reasoning more than factual recall.

The reason this question has not been answered is not lack of interest but lack of a suitable experimental design. Comparing different model families introduces architecture and training confounds. Re-training with varied domain proportions is prohibitively expensive. The gap, at its core, is the absence of \textit{within-family} domain exposure variation that could serve as a covariate in a regression mapping domain composition to benchmark performance.

Our key insight is that this variation already exists --- hidden in the training history of the Pythia model family \cite{Biderman2023Pythia}. Because The Pile stores documents in domain-clustered blocks rather than uniformly shuffling them, different training checkpoints of the same Pythia model have seen different cumulative proportions of Wikipedia, StackExchange, PubMed Abstracts, and other domains. Pythia's identity document index mapping (\texttt{doc\_idx.npy}, confirmed for 134M entries) allows exact reconstruction of which documents each checkpoint consumed --- and therefore how much of each domain it had been exposed to at any point in training. This creates a natural panel dataset: 154 checkpoints across 16 model sizes, each with a distinct domain exposure profile, enabling regression analysis of domain-benchmark specificity without a single new training run.

We operationalize this insight through four sub-hypothesis experiments spanning domain content analysis (h-m1), exposure trajectory extraction (h-e1), Spearman correlation analysis (h-m2), and panel OLS regression (h-m3). The results tell a coherent --- if incomplete --- story. The Pile's domains are measurably distinguishable by automated cognitive task pattern proxies with extremely large effect sizes: entity density distinguishes Wikipedia from Books with $\eta^2 = 0.9915$; narrative coherence distinguishes Books from Wikipedia with $\eta^2 = 0.9142$; formal syntax density identifies GitHub with $\eta^2 = 0.9821$ (all $p \approx 0$; h-m1). Domain exposure trajectories are genuinely non-uniform: 10 of 22 Pile domains show standard deviation exceeding 0.001 across the 154-checkpoint trajectory (h-e1). Together, these two results establish the necessary empirical foundation for domain-benchmark specificity analysis.

The specificity test itself, however, was blocked by a structural data gap: Books3 cumulative exposure is zero for all checkpoints in our 600,000-document first-shard sample, making the Books$\to$HellaSwag prediction (P2) structurally untestable. Preliminary Spearman analysis at the 70M model scale (N=10 checkpoints) shows Wikipedia exposure correlating more strongly with HellaSwag ($\rho = +0.423$) than MMLU ($\rho = -0.391$) --- opposite to our original prediction, though this result reflects known floor effects at 70M scale and unreliable N=10 sampling. The panel OLS regression framework (h-m3) was fully implemented and validated (2,943 lines of code, 21/24 unit tests passing) but could not be executed with the available data.

This paper makes the following contributions:

\textbf{1. Empirical grounding for domain-benchmark specificity analysis.} We provide what is, to our knowledge, the first quantitative measurement showing that The Pile's domains are separable by automated cognitive task pattern proxies with effect sizes exceeding $\eta^2 = 0.91$ --- confirming that domain content differentiation, the assumed mechanism underlying data mixing theories, is empirically observable. (Analysis conducted on domain-representative generated text; $\eta^2$ values represent upper bounds relative to real Pile documents; see Section~\ref{sec:discussion}.)

\textbf{2. Reusable infrastructure for within-family domain exposure analysis.} We extract and validate domain exposure trajectories from Pythia's exact dataloaders using the identity \texttt{doc\_idx} mapping, providing a proven pipeline (\texttt{compute\_domain\_exposure\_trajectories()}, \texttt{build\_domain\_lookup()}, \texttt{step\_to\_sample()}) for future domain attribution studies.

\textbf{3. Identification of a structural data scope requirement.} We document that single-shard Pile samples systematically miss 6 of 22 domains (including Books3 and GitHub), establishing the 134M-document full Pile lookup as the minimum required dataset for domain-benchmark specificity analysis --- a reproducibility warning for future work using partial Pile samples.

\textbf{4. Preliminary directional evidence suggesting scale-dependence.} The 70M reversal (Wikipedia$\to$HellaSwag $>$ Wikipedia$\to$MMLU) is most consistent with a scale floor artifact, but raises the open question of whether domain-benchmark specificity is a scale-emergent phenomenon that requires $\geq$1B parameters to detect.
